\subsection{张量积中元素的秩}

设 E 是右向量空间，F 是同一域 K 上的左向量空间；对每个元素 \(u \in E \otimes_{K} F\) ，在 (29) 下典范地对应一个同态 \(u_{1} \in \operatorname{Hom}_{\mathbb{K}}(\mathbf{E}^{*}, \mathbf{F})\) ；若 \(u = \sum_{i} x_{i} \otimes y_{i}\) ，其中 \(x_{i} \in E, y_{i} \in F\) ，则元素 \(u_{1}\) 是线性映射

\[
x ^ {*} \mapsto \sum_ {i} \langle x ^ {*}, x _ {i} \rangle y _ {i}.\tag{33}
\]

另一方面， \(\mathbf{E} \otimes_{\mathbb{K}} \mathbf{F}\) 被典范地等同于 \(\mathbf{F} \otimes_{\mathbb{K}^0} \mathbf{E}\) ，其中 \(\mathbf{E}\) 被视为左向量空间，而 \(\mathbf{F}\) 被视为反域 \(\mathbf{K}^0\) 上的右向量空间；因此典范地对应于 \(u\) 一个同态 \(u_2 \in \operatorname{Hom}_{\mathbb{K}}(\mathbf{F}^*, \mathbf{E})\) ，它由下式给出

\[
y ^ {*} \mapsto \sum_ {i} x _ {i} \langle y _ {i}, y ^ {*} \rangle ;\tag{34}
\]

\(u_{1}\) （相应地 \(u_{2}\) ）被视为 \(E^{*}\) 到 \(F^{**}\) 中（相应地， \(F^{*}\) 到 \(E^{**}\) 中）的映射时，正是 \(u_{2}\) （相应地 \(u_{1}\) ）的转置。因此 \(u_{1}\) 和 \(u_{2}\) 的秩等于同一个有限数 r，即 F 的子空间 \(u_{1}(E^{*})\) 与 E 的子空间 \(u_{2}(F^{*})\) 的公共维数，其中每一个都典范同构于另一个的对偶（第 5 号）；我们称 r（记为 \(\operatorname{rg}(u)\) ）为 \(E \otimes_{K} F\) 的元素 u 的秩，并称 \(u_{1}(E^{*})\) 和 \(u_{2}(F^{*})\) 为与 u 相关联的子空间（分别属于 F 和 E）。

命题 17. 设 \(u\) 是 \(\mathbf{E} \otimes_{\mathbb{K}} \mathbf{F}\) 的一个元素，而 \(\mathbf{M} \subset \mathbf{E}\) 和 \(\mathbf{N} \subset \mathbf{F}\) 是与其相关联的子空间。对每一个表达式 \(u = \sum_{i=1}^{s} x_i \otimes y_i\) （ \(u\) 的表达式），其中 \(x_i \in \mathbf{E}\) 和 \(y_i \in \mathbf{F}\) （ \(1 \leqslant i \leqslant s\) ），子空间 \(\mathbf{M}\) （相应地 \(\mathbf{N}\) ) 包含在 \(\mathbf{E}\) （相应地 \(\mathbf{F}\) ) 的由诸 \(x_i\) （相应地诸 \(y_i\) ) 生成的子空间中。此外，以下性质是等价的：

\begin{enumerate}
\def\labelenumi{(\alph{enumi})}
\item
  整数 \(s\) 等于 \(u\) 的秩.
\item
  族 \((x_{i})_{1\leqslant i\leqslant s}\) 是M的一组基。
\item
  族 \((y_{i})_{1\leqslant i\leqslant s}\) 是 N 的一组基。
\item
  族 \((x_{i})_{1\leqslant i\leqslant s}\) 和 \((y_{i})_{1\leqslant i\leqslant s}\) 都是自由的。
\end{enumerate}

由 (33)（相应地 (34)）， \(N = u_{1}(E^{*})\) （相应地 \(M = u_{2}(F^{*})\) ) 的每个元素都是诸 \(y_{i}\) （相应地诸 \(x_{i}\) ）的线性组合；由此得第一个断言。若 s = r，则由诸 \(x_{i}\) （相应地 \(y_{i}\) ）生成的、维数 \(\leqslant \dim M\) （相应地 \(\leqslant \dim N\) ）且包含 M（相应地 N）的子空间与它相同，因此 (a) 蕴含 (b) 和 (c)，更不必说 (d)。反之，条件 (b) 和 (c) 中的每一个都依 \(\operatorname{rg}(u)\) 的定义蕴含 (a)。最后，若 (d) 成立，则存在一个族 \((x_{i}^{*})_{1\leqslant i\leqslant s}\) ，其元素属于 \(E^{*}\) ，使得 \(\langle x_{i}, x_{j}^{*}\rangle = \delta_{ij}\) （第 5 号，定理 7 的推论 1），因此由 (33) 可得 \((y_{i})\) 是 N 的一组基，这就完成了证明。

推论 1. \(u\) 的秩是最小整数 \(s\) ，使得存在一个表达式 \(u = \sum_{i=1}^{s} x_i \otimes y_i\) ，其中 \(x_i \in \mathbf{E}\) 和 \(y_i \in \mathbf{F}\) 对于 \(1 \leqslant i \leqslant s\) .

这直接由命题 17 及第 2 号命题 1 得出。

推论 2. 设 K 为一个交换域，E、F 为 K 上的两个向量空间，L 为 K 的一个交换扩域。设 u 为 E ⊗K F 的一个元素，M 和 N 为其相伴子空间，u' 为 u 在 (E ⊗K F)(L) 中的典范像（典范地等同于 E(L) ⊗L F(L)，参见 § 5 第 1 号命题 3）；则 rg(u') = rg(u)，且与 u' 相伴的子空间典范地等同于 M(L) 和 N(L)。

若 \(u = \sum_{i=1}^{r} x_i \otimes y_i\) ，其中族 \((x_i)\) 和 \((y_i)\) 是自由的，则 \(u' = \sum_{i=1}^{r} (1 \otimes x_i) \otimes (1 \otimes y_i)\) 且族 \((1 \otimes x_i)\) 和 \((1 \otimes y_i)\) 分别在 \(E_{(L)}\) 和 \(F_{(L)}\) 中是自由的（§ 5，第 1 号，命题 4）。

